\documentclass[12pt,a4paper]{article}

% --- Настройка языка и шрифтов для LuaLaTeX ---
\usepackage{fontspec}
\defaultfontfeatures{Ligatures=TeX,Scale=MatchLowercase}

\setmainfont{Alice-Regular}[
    Path=font/,
    Extension=.ttf,
    BoldFont = Alice-Regular,
    ItalicFont = Alice-Regular,
    BoldItalicFont = Alice-Regular,
    BoldFeatures = {FakeBold=1.6},
    ItalicFeatures = {FakeSlant=0.2},
    BoldItalicFeatures = {FakeBold=1.6,FakeSlant=0.2}
]

\newfontfamily\cyrillicfont{Alice-Regular}[
    Path=font/,
    Extension=.ttf,
    BoldFont = Alice-Regular,
    ItalicFont = Alice-Regular,
    BoldItalicFont = Alice-Regular,
    BoldFeatures = {FakeBold=1.6},
    ItalicFeatures = {FakeSlant=0.2},
    BoldItalicFeatures = {FakeBold=1.6,FakeSlant=0.2}
]

\usepackage{polyglossia}
\setmainlanguage{russian}
\setsansfont{TeX Gyre Heros}
\setmonofont{TeX Gyre Cursor}

\usepackage{graphicx}        					          % Графика
\usepackage[protrusion=false]{microtype}       	% Микротипографика
\usepackage{amsmath, amsfonts, amssymb, amsthm}	% Математика
\usepackage{braket} 							              % Braket нотация для квантовой механики
\usepackage{unicode-math} 						          % Современный пакет для математики в LuaLaTeX
\usepackage{longtable}       					          % Таблицы с разрывом страниц
\usepackage{tikz}           					          % Рисование
\usepackage[europeanresistors]{circuitikz}		  % Рисование электрических схем в европейском стиле
\usepackage{listings}							              % Для отображения кода
\usepackage{tabularx}       					          % Для таблиц во всю ширину
\usepackage{array}		  						            % Для расширенных возможностей таблиц
\usepackage{float}          					          % Для точного позиционирования фигур
\usepackage{pgfplots}                           % Для построения графиков
\usepackage{caption}                            % Для настройки подписей к рисункам и таблицам
\usepackage{xfp}                                % Для вычислений в LaTeX
\pgfplotsset{compat=1.18}
\usepackage[a4paper,left=20mm,right=20mm,top=20mm,bottom=20mm]{geometry}

\usetikzlibrary{arrows.meta}

% --- Колонтитулы ---
\usepackage{fancyhdr}
\pagestyle{fancy}
\fancyhf{}

\renewcommand{\headrulewidth}{0.4pt}
\renewcommand{\footrulewidth}{0.4pt}
\renewcommand\lstlistingname{Исходный код}

\lstdefinestyle{main}{
	backgroundcolor=\color{gray!8},   
	commentstyle=\color{green!50!black},
	keywordstyle=\color{blue!80!black},
	numberstyle=\small\color{darkgray},
	stringstyle=\rmfamily\color{orange!90!black},
	basicstyle=\ttfamily\small,
	breakatwhitespace=false,         
	breaklines=true,                   
	keepspaces=true,                 
	numbers=left,                    
	numbersep=8pt,                  
	showspaces=false,                
	showstringspaces=false,
	showtabs=false,                  
	tabsize=4
}	

\lstset{style=main}

\newcolumntype{C}{>{\centering\arraybackslash}X}
\renewcommand{\arraystretch}{1.25}

% =========================================================
%         СЛУЖЕБНЫЕ КОМАНДЫ ДЛЯ ВЫЧИСЛЕНИЙ
% =========================================================
\newcommand{\CalcZero} [1]{\fpeval{round((#1),0)}}
\newcommand{\CalcOne}  [1]{\fpeval{round((#1),1)}}
\newcommand{\CalcTwo}  [1]{\fpeval{round((#1),2)}}
\newcommand{\CalcThree}[1]{\fpeval{round((#1),3)}}
\newcommand{\CalcFour} [1]{\fpeval{round((#1),4)}}
\newcommand{\CalcFive} [1]{\fpeval{round((#1),5)}}

% ---------------------------------------------------------
% ОБЩИЕ ДАННЫЕ
% ---------------------------------------------------------
\def\VariantN{1}          					      % Номер варианта
\def\StudentName{Андреев Иван Алексеевич} % ФИО студента
\def\LabNumber{14}							          % Номер лабораторной работы
\def\Discipline{Атомная физика}				    % Название дисциплины
\def\LabTitle{Эффект Холла}	              % Название лабораторной работы

\title{
	МИНОБРНАУКИ РОССИИ

	федеральное государственное автономное образовательное учреждение

	высшего образования

	«Национальный исследовательский университет
	«Московский институт электронной техники»
	\vspace{2em}
}
\author{
}
\date{}

\begin{document}

\maketitle
\thispagestyle{fancy}

\begin{center}
\Large Отчёт по Лабораторной работе $\LabNumber$\\
по дисциплине «\Discipline»\\
«\LabTitle»
\end{center}

\begin{center}
\Large Вариант \VariantN
\end{center}

\vspace{12em}

\cfoot{Москва 2026}

\pagebreak

\fancyhead[L]{\LabTitle} 
\fancyhead[R]{\Discipline}
\fancyfoot[C]{\thepage}

\section{Зависимость напряжения Холла от тока через образец}

\def\B{\fpeval{0.2}} % Магнитная индукция в Теслах
\def\b{\fpeval{1.3 * 10^{-3}}}  % Толщина образца в мм
\def\gammaI{\fpeval{1.6}} % Коэффициент пропорциональности для эффекта Холла

Дано
\begin{equation*}
  U_H(I_p) \quad  B = \B \text{ Тл} \quad d = 10 \text{ мм} \quad l = 20 \text{ мм} \quad b = \b \text{ мм}
\end{equation*}

\begin{equation*}
  U_H = \frac{I_pB}{b}R
\end{equation*}

\def\R1{\fpeval{(\gammaI * \b) / (\B)}} % Сопротивление образца, найденное из первого эксперимента

\begin{equation*}
	\begin{aligned}
		\gamma = \frac{R_\text{х} B}{b} \iff R_\text{х} = \frac{\gamma b}{B} \quad &   \\
		& R_\text{х} = \frac{\gammaI \cdot \b}{\B} = \R1
	\end{aligned}
\end{equation*}

Найти $R$

\section{Зависимость напряжения Холла от магнитного поля}

\def\Ip{\fpeval{30 * 10^{-3}}} % Ток через образец в Амперах
\def\gammaB{\fpeval{243.9 * 10^(-3)}} % Коэффициент пропорциональности для эффекта Холла

\def\R2{\fpeval{(\gammaB * \b) / (\Ip)}} % Сопротивление образца, найденное из второго эксперимента
\def\p{\fpeval{1 / (1.6 * 10^(-19) * \R2)}} % Концентрация носителей заряда, найденная из второго эксперимента

Дано
\begin{equation*}
  U_H(B) \quad  I_p = \fpeval{\Ip * 10^3} \text{ мА} \quad d = 10 \text{ мм} \quad l = 20 \text{ мм} \quad b = \b \text{ мм}
\end{equation*}

\begin{equation*}
  U_H = \frac{I_pB}{b}R \quad R = \frac{1}{ep}
\end{equation*}

\begin{equation*}
	\begin{aligned}
		\gamma = \frac{R_\text{х} I_p}{b} \iff R_\text{х} = \frac{\gamma b}{B} \quad &   \\
		& R_\text{х} = \frac{\gammaB \cdot \b}{\Ip} = \R2
	\end{aligned}
\end{equation*}

\begin{equation*}
	\begin{aligned}
		R_\text{х} = \frac{1}{ep} \iff p = \frac{1}{eR_\text{х}} \quad &   \\
		& p = \CalcFive{\p * 10^{-20}} \cdot 10^{20} \text{ м$^{-3}$}
	\end{aligned}
\end{equation*}

Найти $R$, $p$
 
\end{document}